% GENERATED FILE. Edit canonical lessons, then run npm run generate:books.
% canonical-source-hash: fnv1a64:eb6903988a97052f
% canonical-lessons: PT-C110-enfermeira, PT-C110-cozinheiro, PT-C110-empregado, PT-C110-agricultor, PT-C110-padeiro

\chapter{Nurse, Cook, Waiter, Farmer, Baker}
\label{ch:pt-nurse-cook-waiter-farmer-baker}
\hlchaptermodality{110}

\begin{chapteropening}
\textbf{By the end of this chapter:} \emph{I can say a nurse, a cook, a waiter, a farmer and a baker.}

Complete the last lesson of chapter 110: I can say a nurse, a cook, a waiter, a farmer and a baker.
\end{chapteropening}

\section[a enfermeira]{a enfermeira — a nurse}

\label{lesson:PT-C110-enfermeira}



\begin{quote}
Before the new one: say the Portuguese for pepper, then the Portuguese for a planet.
\end{quote}

\subsection*{You'll want to know: a enfermeira}
\textbf{a enfermeira} — \textquotedblleft{}a nurse\textquotedblright{}.

\begin{cousinweb}
\textbf{enfermeira} is built from an old word for \textquotedblleft{}sick\textquotedblright{} and \emph{-eira}. That word $\leftarrow$ Latin \emph{infirmus}, \textquotedblleft{}weak, not strong\textquotedblright{} (EtymDB 2.1).
\end{cousinweb}

\begin{grammarlens}[title={What you've built}]
One more word for this chapter.
\end{grammarlens}

\subsection*{Your turn}
\begin{itemize}
\raggedright
  \item \emph{Say it:} \emph{a enfermeira}
  \item \emph{Say it:} \emph{a enfermeira}, once more
  \item \emph{Say it:} say \emph{a enfermeira}
  \item \emph{Recall:} say the Portuguese for pepper, then the Portuguese for a planet, then say \emph{a enfermeira} again
\end{itemize}

\begin{tcolorbox}[breakable,colback=teal!4,colframe=teal!35!black,title={Before you move on}]
What does a enfermeira mean? (\textquotedblleft{}A nurse\textquotedblright{}.) Say it once more.
\end{tcolorbox}
\section[o cozinheiro]{o cozinheiro — a cook}

\label{lesson:PT-C110-cozinheiro}



\begin{quote}
Before the new one: say the Portuguese for a planet, then the Portuguese for a nurse.
\end{quote}

\subsection*{You'll want to know: o cozinheiro}
\textbf{o cozinheiro} — \textquotedblleft{}a cook\textquotedblright{}.

\begin{cousinweb}
\textbf{cozinheiro} is built from \emph{cozinha}, \textquotedblleft{}kitchen\textquotedblright{}, and \emph{-eiro}, \textquotedblleft{}one who works at\textquotedblright{}.
\end{cousinweb}

\begin{grammarlens}[title={What you've built}]
One more word for this chapter.
\end{grammarlens}

\subsection*{Your turn}
\begin{itemize}
\raggedright
  \item \emph{Say it:} \emph{o cozinheiro}
  \item \emph{Say it:} \emph{o cozinheiro}, once more
  \item \emph{Say it:} say \emph{o cozinheiro}
  \item \emph{Recall:} say the Portuguese for a planet, then the Portuguese for a nurse, then say \emph{o cozinheiro} again
\end{itemize}

\begin{tcolorbox}[breakable,colback=teal!4,colframe=teal!35!black,title={Before you move on}]
What does o cozinheiro mean? (\textquotedblleft{}A cook\textquotedblright{}.) Say it once more.
\end{tcolorbox}
\section[o empregado]{o empregado — a waiter; an employee}

\label{lesson:PT-C110-empregado}



\begin{quote}
Before the new one: say the Portuguese for a nurse, then the Portuguese for a cook.
\end{quote}

\subsection*{You'll want to know: o empregado}
\textbf{o empregado} — \textquotedblleft{}a waiter; an employee\textquotedblright{}. In Brazil, \textbf{o garçom, for a waiter}.

\begin{cousinweb}
\textbf{empregado} is \textquotedblleft{}employed\textquotedblright{}: the participle of the verb \textquotedblleft{}to employ\textquotedblright{}, which $\leftarrow$ Latin \emph{implicare}, \textquotedblleft{}to involve, to enfold\textquotedblright{} (Lewis \& Short).
\end{cousinweb}

\begin{grammarlens}[title={What you've built}]
One more word for this chapter.
\end{grammarlens}

\subsection*{Your turn}
\begin{itemize}
\raggedright
  \item \emph{Say it:} \emph{o empregado}
  \item \emph{Say it:} \emph{o empregado}, once more
  \item \emph{Say it:} say \emph{o empregado}
  \item \emph{Recall:} say the Portuguese for a nurse, then the Portuguese for a cook, then say \emph{o empregado} again
\end{itemize}

\begin{tcolorbox}[breakable,colback=teal!4,colframe=teal!35!black,title={Before you move on}]
What does o empregado mean? (\textquotedblleft{}A waiter; an employee\textquotedblright{}.) Say it once more.
\end{tcolorbox}
\section[o agricultor]{o agricultor — a farmer}

\label{lesson:PT-C110-agricultor}



\begin{quote}
Before the new one: say the Portuguese for a cook, then the Portuguese for a waiter; an employee.
\end{quote}

\subsection*{You'll want to know: o agricultor}
\textbf{o agricultor} — \textquotedblleft{}a farmer\textquotedblright{}.

\begin{cousinweb}
\textbf{agricultor} $\leftarrow$ Latin \emph{agricultor}, \textquotedblleft{}field-tiller\textquotedblright{} (EtymDB 2.1), from \emph{ager}, \textquotedblleft{}field\textquotedblright{}, and \emph{colere}, \textquotedblleft{}to till\textquotedblright{} — a learned borrowing from Latin.
\end{cousinweb}

\begin{grammarlens}[title={What you've built}]
One more word for this chapter.
\end{grammarlens}

\subsection*{Your turn}
\begin{itemize}
\raggedright
  \item \emph{Say it:} \emph{o agricultor}
  \item \emph{Say it:} \emph{o agricultor}, once more
  \item \emph{Say it:} say \emph{o agricultor}
  \item \emph{Recall:} say the Portuguese for a cook, then the Portuguese for a waiter; an employee, then say \emph{o agricultor} again
\end{itemize}

\begin{tcolorbox}[breakable,colback=teal!4,colframe=teal!35!black,title={Before you move on}]
What does o agricultor mean? (\textquotedblleft{}A farmer\textquotedblright{}.) Say it once more.
\end{tcolorbox}
\section[o padeiro]{o padeiro — a baker}

\label{lesson:PT-C110-padeiro}



\begin{quote}
Before the new one: say the Portuguese for a waiter; an employee, then the Portuguese for a farmer.
\end{quote}

\subsection*{You'll want to know: o padeiro}
\textbf{o padeiro} — \textquotedblleft{}a baker\textquotedblright{}.

\begin{cousinweb}
\textbf{padeiro} $\leftarrow$ Vulgar Latin \emph{*panatarius} (EtymDB 2.1), from Latin \emph{panis}, \textquotedblleft{}bread\textquotedblright{} — inherited.
\end{cousinweb}

\begin{grammarlens}[title={What you've built}]
One more word for this chapter.
\end{grammarlens}

\subsection*{Your turn}
\begin{itemize}
\raggedright
  \item \emph{Say it:} \emph{o padeiro}
  \item \emph{Say it:} \emph{o padeiro}, once more
  \item \emph{Say it:} say \emph{o padeiro}
  \item \emph{Recall:} say the Portuguese for a waiter; an employee, then the Portuguese for a farmer, then say \emph{o padeiro} again
\end{itemize}

\begin{tcolorbox}[breakable,colback=teal!4,colframe=teal!35!black,title={Before you move on}]
What does o padeiro mean? (\textquotedblleft{}A baker\textquotedblright{}.) Say it once more.
\end{tcolorbox}
